\chapter{Introduction to File Handling in Python}

\section{Introduction: Real-Life Hook}

Students deal with information stored in files every day:  

\begin{itemize}
    \item Saving homework or notes  
    \item Keeping contact lists  
    \item Storing exam scores  
\end{itemize}

Python allows programs to \textbf{read from and write to files}, so data can be stored permanently and reused later.

\section{Concept: What is a File?}

\begin{definition}{File}
A \textbf{file} is a collection of information stored on a computer, which can be text or binary data.  
Python provides built-in functions to read and write text files.
\end{definition}

\section{Opening and Closing Files}

\begin{verbatim}
# Open a file in write mode
f = open("data.txt", "w")  # 'w' = write mode
f.write("Hello, Python!\n")
f.close()
\end{verbatim}

- \texttt{"w"} = write mode (creates new file or overwrites existing)  
- \texttt{"r"} = read mode  
- \texttt{"a"} = append mode (adds to existing file)  
- Always \textbf{close} files to save changes  

\subsection{Using with Statement (Recommended)}

\begin{verbatim}
with open("data.txt", "w") as f:
    f.write("Hello, Python!\n")
# File automatically closes
\end{verbatim}

\section{Reading from Files}

\begin{verbatim}
with open("data.txt", "r") as f:
    content = f.read()
    print(content)
\end{verbatim}

- \texttt{read()} reads the entire file  
- \texttt{readline()} reads one line at a time  
- \texttt{readlines()} returns all lines as a list  

\subsection{Example: Reading Line by Line}

\begin{verbatim}
with open("data.txt", "r") as f:
    for line in f:
        print(line.strip())  # Remove newline characters
\end{verbatim}

\section{Appending to Files}

\begin{verbatim}
with open("data.txt", "a") as f:
    f.write("New line added\n")
\end{verbatim}

\section{Mini-Project: Contact List Program}

Create a program to save and read contact names:

\begin{verbatim}
# Write contacts
contacts = ["Amina", "Rafi", "Sadia"]

with open("contacts.txt", "w") as f:
    for contact in contacts:
        f.write(contact + "\n")

# Read contacts
with open("contacts.txt", "r") as f:
    print("Contact List:")
    for line in f:
        print(line.strip())
\end{verbatim}

Output:

\begin{verbatim}
Contact List:
Amina
Rafi
Sadia
\end{verbatim}

\section{Step-by-Step Diagram: File Handling Flow}

\begin{center}
\begin{tikzpicture}[node distance=1.5cm]
\node (start) [draw, ellipse, fill=green!20] {Start Program};
\node (open) [draw, rectangle, fill=yellow!20, below of=start] {Open File (r/w/a)};
\node (process) [draw, rectangle, fill=orange!20, below of=open] {Read/Write Data};
\node (close) [draw, rectangle, fill=yellow!20, below of=process] {Close File};
\node (end) [draw, ellipse, fill=green!20, below of=close] {End Program};
\draw[->] (start) -- (open);
\draw[->] (open) -- (process);
\draw[->] (process) -- (close);
\draw[->] (close) -- (end);
\end{tikzpicture}
\end{center}

\section{Key Concepts}

\begin{itemize}
    \item Files store data permanently on a computer  
    \item Python can \textbf{read} and \textbf{write} files using built-in functions  
    \item Always \textbf{close files} after use or use \texttt{with} statement  
    \item File modes: \texttt{"r"} = read, \texttt{"w"} = write, \texttt{"a"} = append  
    \item Use loops to read files line by line  
\end{itemize}

\section{End-of-Chapter Exercises}

\begin{exercise}{Very Short Questions}
\begin{enumerate}
    \item What is a file in Python?  
    \item Name three file modes in Python.  
    \item Why should files be closed after writing?  
\end{enumerate}
\end{exercise}

\begin{exercise}{Short Answer Questions}
\begin{enumerate}
    \item Write Python code to read a file and print its content.  
    \item How is \texttt{with open()} different from \texttt{open()}?  
    \item Explain the difference between \texttt{read()}, \texttt{readline()}, and \texttt{readlines()}.  
\end{enumerate}
\end{exercise}

\begin{exercise}{Activity-Based Questions}
\begin{enumerate}
    \item Write a program to take 5 names as input and save them in a file. Then read the file and display all names.  
    \item Create a program to append a new score to an existing file of student marks.  
    \item Read a text file and count the total number of words in it.  
\end{enumerate}
\end{exercise}

\section{Teacher Notes}

\marginnote{
\textbf{Learning Objectives:}  
\begin{itemize}
    \item Students understand how to open, read, write, and append files  
    \item Students can use loops to process file data line by line  
    \item Students learn to use \texttt{with open()} for safe file handling  
\end{itemize}

\textbf{Teaching Tips:}  
\begin{itemize}
    \item Demonstrate writing to a file, reading it, and appending data  
    \item Relate examples to daily life: homework, contact lists, scores  
\end{itemize}

\textbf{Assessment:}  
\begin{itemize}
    \item Check correct reading and writing of files  
    \item Observe use of loops with file data  
\end{itemize}
}